\section{The first higher images through rank five}
\label{sch:sec:first-higher}

The first higher scalar coefficient for each of $f_4$ and $f_5$
is one-dimensional. Complete-flag line classes determine its binary
and three-constituent Hall images.
All ranks in this section satisfy $2\leq n\leq5$.
The scalar points, coefficients and orientations remain those of
Sections~\ref{sch:sec:scalar}--\ref{sch:sec:evaluation}.

\begin{lemma}[Two adjacent scalar degrees]
\label{sch:lem:next-degrees}
At a scalar atlas point,
\[
 H^{-2n+1}(P_{n,0})=0,\qquad
 H^{-2n+2}(P_{n,0})=\mathbb Q.
\]
\end{lemma}

\begin{proof}
Set $N=n^2+n=\dim_{\mathbb C}\mathcal C_n$.
Remove the pairs for which $A$ has at most $n-2$ distinct eigenvalues.
Their complex dimension is at most $N-2$, by the Jordan calculation.
This removal preserves Borel--Moore homology in degrees $2N-1$
and $2N-2$. For $n=2$, the removed locus is empty.
Let $O$ be the distinct-eigenvalue locus in the remaining space $X$.
Its ordered cover has, up to a contractible affine factor, the form
\[
 \operatorname{GL}_n(\mathbb C)/T_n
                    \times\operatorname{Conf}_n(\mathbb C).
\]
Here $T_n$ is the diagonal torus. The first factor has $H^1=0$.
Its $H^2$ consists of the $n$
Chern roots modulo their sum. Hence its symmetric-group invariants
in degree two vanish.

The degree-one cohomology of ordered configurations has basis
$a_{ij}=(2\pi i)^{-1}d\log(z_i-z_j)$, for $i<j$, and its
cohomology in degree two is generated by their products.
Here is a direct induction for the required generation statement.
Forgetting the last point is a locally trivial bundle with fibre
$\mathbb C\setminus\{z_1,\ldots,z_{n-1}\}$.
Local trivializations follow by moving the points in disjoint disks
and extending those motions to the plane with compact support.
The classes $a_{1n},\ldots,a_{n-1,n}$ restrict to the basis dual
to the small positive circles in the punctured-plane fibre.
These global classes supply permanent cycles for the whole fibre
cohomology. The Serre spectral sequence therefore collapses. The cohomology is a
free module over the base cohomology on $1,a_{1n},\ldots,a_{n-1,n}$.
Induction from $n=1$ proves the assertion, including the stated
degree-one basis.

The symmetric group permutes the unordered-edge basis $a_{ij}$.
It has a one-dimensional invariant subspace in degree one.
The exterior square of the full degree-one representation has no invariants.
Indeed, for two distinct edges sharing a vertex, the transposition
of the other endpoints interchanges the two edges.
For disjoint edges, interchanging their two pairs of endpoints has
the same effect. Each operation negates the corresponding wedge.
Its average over the symmetric group is therefore zero.
Since invariants over $\mathbb Q$ preserve surjections, the quotient
$H^2(\operatorname{Conf}_n)$ also has no invariants.
Consequently $H^1(O)=\mathbb Q$ and $H^2(O)=0$.

The complement $X\setminus O$ has two components of complex
dimension $N-1$. They have one double eigenvalue of $A$, with
respectively a Jordan block and a semisimple double eigenspace.
All remaining eigenvalues are distinct. Centralizer bundles show
that both strata are irreducible. Their closures are distinct by
the same commuting-eigenvalue and minimal-polynomial conditions used
in Lemma~\ref{sch:lem:rank-three-bm}.
The discriminant has multiplicities one and two on these components.
To check this, use the two-by-two blocks
\[
 \begin{pmatrix}0&1\\t&0\end{pmatrix},\qquad
 \operatorname{diag}(t,-t),
\]
and append fixed, pairwise distinct eigenvalues away from zero.
All additional factors in the discriminant are units near $t=0$.
The cyclic-$A$ chart in the first case and cyclic-$C$ chart in the
second give smooth transverse coordinates, as before.

The closed-open sequence thus has the exact segment
\[
 0\longrightarrow H^{\mathrm{BM}}_{2N-1}(X)
 \longrightarrow\mathbb Q\xrightarrow{\ (1,2)\ }\mathbb Q^2
 \longrightarrow H^{\mathrm{BM}}_{2N-2}(X)\longrightarrow0.
\]
The first group is zero and the last is one-dimensional.
Restoring the removed locus preserves these groups.
The scalar support isomorphism \eqref{sch:eq:scalar-bm} gives the
two asserted degrees.
\end{proof}

For the complete flag in dimension $n$, let
$\kappa_{n,k}:\mathbb Q_{\mathfrak M_n}[2n-2]\to\mathcal P_n$
be the complete-flag morphism with $c_1(L_k)$ inserted.
Here $L_k$ is its $k$th successive quotient line.
Thus $\kappa_{3,k}=\kappa_k$ and
$\kappa_{2,2}=\tau_2$.

\begin{proposition}[Normalized first higher classes]
\label{sch:prop:higher-generators}
The morphism
\[
 \tau_n=\frac{\kappa_{n,n}}{(n-1)!}
       :\mathbb Q_{\mathfrak M_n}[2n-2]\longrightarrow\mathcal P_n
\]
has a nonzero scalar class $w_n$ in degree $-2n+2$.
It generates the group of Lemma~\ref{sch:lem:next-degrees}.
For each $1\leq k\leq n$, its scalar line decorations satisfy
\[
 \kappa_{n,k}(1)=(2k-n-1)(n-2)!\,w_n.
\]
These generators agree with $w_2$ and $w_3$ already constructed.
\end{proposition}

\begin{proof}
Take a critical triple with one scalar plane and $n-2$ separated
scalar lines. The spectral projectors split its local coefficient
into the rank-two coefficient and the line coefficients, together
with the nondegenerate cross-block normal quadratic.
The same invertible commutator operator and complex rescaling as in
Lemma~\ref{sch:lem:three-decorations} give this decomposition.
There are $p_n-1$ cross eigenline pairs. Their $B$-polarized sign
$(-1)^{p_n-1}$ is canceled by the generator-conversion sign
$(-1)^{1-p_n}$. Thus every flag component contributes $\nu_2$
with coefficient $+1$.

Choose the two positions of the lines in the double eigenspace.
The other $n-2$ labeled lines can occupy the remaining positions
in $(n-2)!$ orders. Each such choice gives one $\mathbb P^1$
component of the flag fibre. On it, $c_1(L_k)$ is $-h$ if
position $k$ is the first double-eigenspace line, $h$ if it is
the second, and zero otherwise.
There are $(n-k)(n-2)!$ negative contributions and
$(k-1)(n-2)!$ positive contributions.
The value of $\kappa_{n,k}$ is therefore
$(2k-n-1)(n-2)!\,w_2$ on this stratum, including the line factors.

For $k=n$ this is $(n-1)!w_2\ne0$.
Positive-scaling equivariance of the actual cohomology section forces
its scalar origin germ to be nonzero.
The scalar group is one-dimensional by
Lemma~\ref{sch:lem:next-degrees}.
Subtracting the claimed multiple of $\tau_n$ from each decoration
gives a section whose scalar germ is some multiple of $w_n$.
Its value on the separated-plane stratum determines this multiple,
by the same zero-germ scaling argument.
This proves the formula and the normalizations.
\end{proof}

Fix $d+e=n\leq5$ and let $D_{d,e}$ be the actual scalar incoming
complex for $\mu_{d,e}$. Let $U\subset\mathbb C^n$ be the
tautological $d$-plane on $\operatorname{Gr}(d,n)$ and put
$H=c_1(\mathbb C^n/U)$.
Define the following scalar classes in degree $-2n+2$:
\[
 \begin{aligned}
 A_{d,e}&=H\,(v_d\otimes v_e),\\
 B_{d,e}&=\bigl(Rb_*a^*(\tau_d\boxtimes\sigma_e)
                             \circ\eta[2n-2]\bigr)_0(1)
                                      &&(d\geq2),\\
 C_{d,e}&=\bigl(Rb_*a^*(\sigma_d\boxtimes\tau_e)
                             \circ\eta[2n-2]\bigr)_0(1)
                                      &&(e\geq2).
 \end{aligned}
\]
The maps $a,b,\eta$ refer to the same full critical correspondence
and its adjunction unit as in \eqref{sch:eq:divided}.
The letters $A_{d,e},B_{d,e},C_{d,e}$ denote coefficient classes,
not representation matrices.

\begin{theorem}[Binary first higher images]
\label{sch:thm:binary-higher}
The displayed classes form a basis of
$H^{-2n+2}(D_{d,e})$, whose dimension is
$1+\mathbf 1_{d\geq2}+\mathbf 1_{e\geq2}$.
The actual Hall map has values
\[
 \begin{aligned}
 \mu_{d,e}(A_{d,e})
    &=\frac{(n-2)!}{(d-1)!(e-1)!}\,w_n,\\
 \mu_{d,e}(B_{d,e})
    &=\frac{(d-e-1)(n-2)!}{(d-1)!e!}\,w_n &&(d\geq2),\\
 \mu_{d,e}(C_{d,e})
    &=\frac{(n-1)!}{d!(e-1)!}\,w_n &&(e\geq2).
 \end{aligned}
\]
In particular this map is surjective and has rank one.
\end{theorem}

\begin{proof}
The incoming coefficient has no group one degree above its lowest
degree, by Lemma~\ref{sch:lem:next-degrees} and $P_1=\mathbb Q[2]$.
Its groups two degrees above the lowest come from the first higher
class of either nontrivial-rank factor.
The Grassmannian has no odd cohomology and has one-dimensional $H^2$.
Thus the only terms of total degree $-2n+2$ are its $H^2$ with
the lowest coefficient and its $H^0$ with these higher coefficients.
The same differential argument as in
Proposition~\ref{sch:prop:degree-four} makes all these terms permanent.
The classes $B_{d,e},C_{d,e}$ lift their separate nonzero edge classes,
whereas $A_{d,e}$ generates the flag contribution. This proves the basis.

Refine the partial flag to a complete flag and commute Chern
insertions through the proper support comparison.
The three incoming classes become respectively
\[
 \frac{\sum_{k=d+1}^{n}c_1(L_k)}{d!e!},\qquad
 \frac{c_1(L_d)}{(d-1)!e!},\qquad
 \frac{c_1(L_n)}{d!(e-1)!}
\]
under the complete-flag coefficient map.
Proposition~\ref{sch:prop:higher-generators} evaluates them.
For the first numerator,
$\sum_{k=d+1}^{n}(2k-n-1)=de$.
The other two numerators are $d-e-1$ and $n-1$.
These identities give the stated coefficients.
The first coefficient is a positive integer, proving surjectivity.
\end{proof}

For the three target ranks the resulting matrices, in the displayed
geometric bases and the normalized $w_n$, are
\[
\begin{array}{c|c|c}
 (d,e)&\text{degree}&\text{matrix on }(A,B,C)\text{ with absent terms omitted}\\\hline
 (2,1)&-4&(1,0)\\
 (2,2)&-6&(2,-1,3)\\
 (3,1)&-6&(1,1)\\
 (4,1)&-8&(1,2)\\
 (2,3)&-8&(3,-2,6).
\end{array}
\]

The same Chern insertions determine the three-constituent comparison.
For a scalar flag with dimensions $(d,e,f)$, set
$X=c_1(\mathbb C^n/E_d)$ and
$Y=c_1(\mathbb C^n/E_{d+e})$.
Their products with the lowest coefficient give the two flag classes.
Include the actual lifts with $\tau$ in each constituent of rank
at least two. For $(1,1,1),(2,1,1),(2,2,1)$, these form bases by
the preceding spectral-sequence argument on the partial flag variety.
Its Schubert cells give two independent degree-two classes $X,Y$
and no odd cohomology.
The common derived flag map has matrices
\[
\begin{array}{c|c|c|c}
 (d,e,f)&\text{degree}&\text{basis order}&\text{image coefficients in }w_n\\\hline
 (1,1,1)&-4&(X,Y)&(2,2)\\
 (2,1,1)&-6&(X,Y,\tau_2\text{ in the first factor})&(4,3,-2)\\
 (2,2,1)&-8&(X,Y,\tau_2\text{ in the first, second factors})&(9,6,-6,6).
\end{array}
\]
Indeed, complete-flag refinement divides the two flag insertions
by $d!e!f!$ and replaces them by the two corresponding suffix sums
of the $c_1(L_k)$. A constituent insertion divides instead by the
factorial one smaller in that constituent and selects its last line.
The values follow from Proposition~\ref{sch:prop:higher-generators}.
The two parenthesizations give these same maps because the ordered
support products, Chern pullbacks and complex Thom integrations commute.
This determines the first higher scalar images and their stated
three-constituent compatibility through rank five.
